\section{Reproducibility Appendix and Artifact Manifest}

\subsection{Environment and Execution Protocol}

All experiments and report figures were generated in the project-local virtual environment:
\begin{lstlisting}[basicstyle=\ttfamily\footnotesize,frame=single]
source /Users/tahamajs/Documents/uni/venv/bin/activate
cd /Users/tahamajs/Documents/uni/NLP/nlp-assignments-spring-2023/NLP_UT/CA2
pip install tldextract
jupyter nbconvert --to notebook --execute answer/NLP_HW2.ipynb \
  --output /tmp/ca2_exec_check.ipynb --ExecutePreprocessor.timeout=3600
\end{lstlisting}

The executed notebook was then promoted back to the deliverable path so that outputs remain embedded in the submission artifact.

\subsection{Deterministic Data Contracts}

The notebook now uses fixed project paths, removing ambiguity from execution context:
\begin{itemize}
    \item \texttt{datasets/q1/emails\_1.csv}
    \item \texttt{datasets/q1/emails\_2.csv}
    \item \texttt{datasets/q2/urls.csv}
\end{itemize}

Label normalization is explicit and validated at runtime:
\begin{itemize}
    \item \texttt{spam/ham} $\rightarrow 1/0$ for \texttt{emails\_1.csv}
    \item \texttt{Prediction} (or \texttt{status}) coerced to integer for \texttt{emails\_2.csv}
    \item \texttt{phishing/legitimate} $\rightarrow 1/0$ for \texttt{urls.csv}
\end{itemize}

This design prevents silent fallback behavior and ensures reproducible class semantics across runs.

\subsection{Figure Manifest}

Table~\ref{tab:manifest} lists all figures that are generated automatically by notebook execution and consumed by this report.

\begin{table*}[t]
\centering
\caption{Generated figure artifact manifest}
\label{tab:manifest}
\begin{tabular}{p{4.2cm}p{8.8cm}p{3.2cm}}
\toprule
Filename & Purpose & Used in section \\
\midrule
\texttt{q1\_confusion\_matrices.png} & Side-by-side confusion matrices for from-scratch Q1 models & Q1 results \\
\texttt{q1\_roc\_pr\_curves\_lr.png} & ROC and Precision-Recall curves for Q1 Logistic Regression & Q1 results \\
\texttt{q1\_top\_words\_spam\_ham.png} & Top class-specific BoW terms from \texttt{emails\_1.csv} & Q1 results \\
\texttt{q1\_model\_comparison\_bar.png} & Metric bar chart for Q1 model comparison & Q1 results \\
\texttt{q2\_structural\_distributions.png} & Histograms of structural URL features & Q2 results \\
\texttt{q2\_correlation\_heatmap\_all\_features.png} & Correlation structure of combined URL features & Q2 results \\
\texttt{q2\_structural\_boxplots.png} & Class-conditional spread for key structural features & Q2 results \\
\texttt{q2\_pairplot\_sample.png} & Pairwise sampled visualization of combined features & Q2 results \\
\texttt{q2\_top\_mean\_abs\_correlations.png} & Top features by mean absolute correlation & Q2 results \\
\texttt{q2\_roc\_combined\_features.png} & ROC comparison on combined feature set & Q2 results \\
\texttt{q2\_pr\_combined\_lr.png} & PR curve for combined-feature Logistic Regression & Q2 results \\
\texttt{q2\_lr\_feature\_importance.png} & Positive/negative coefficient importance plot & Q2 results \\
\bottomrule
\end{tabular}
\end{table*}

\subsection{Complete Feature Dictionary for Q2}

For reproducibility and future extension, Table~\ref{tab:feature-dict} documents all URL features used in the combined model, grouped by engineering family.

\begin{table*}[t]
\centering
\caption{Complete feature dictionary used in the combined Q2 model}
\label{tab:feature-dict}
\begin{tabular}{p{3.0cm}p{2.2cm}p{9.2cm}}
\toprule
Feature & Family & Description and intended signal \\
\midrule
\texttt{host\_len} & Structural & Length of host segment; unusually long hosts can indicate obfuscation or generated domains. \\
\texttt{path\_len} & Structural & Number of characters in URL path; phishing pages often include long nested paths. \\
\texttt{url\_len} & Structural & Total URL length; extreme length can correlate with attack payload stuffing. \\
\texttt{num\_params} & Structural & Number of query parameters; excessive parameters may indicate tracking or camouflage. \\
\texttt{has\_fragment} & Structural & Binary indicator for fragment usage; sometimes used to hide payload-like tokens. \\
\texttt{digits\_in\_host} & Structural & Count of numeric characters in host; high counts can indicate non-human domain patterns. \\
\texttt{length\_url} & Statistical & Raw URL length from content-oriented extraction; provides a second length perspective. \\
\texttt{ratio\_digits\_url} & Statistical & Ratio of digits to full URL length; captures aggressive numeric token usage. \\
\texttt{longest\_words\_raw} & Statistical & Length of longest alphabetic token after split; can expose generated token artifacts. \\
\texttt{sub\_len} & Creative & Subdomain length; deeply nested or very long subdomains are common in phishing. \\
\texttt{has\_ip} & Creative & Binary indicator for IP-literal host usage instead of domain names. \\
\texttt{at\_sign} & Creative & Presence of '@' in URL; classic obfuscation pattern in malicious links. \\
\texttt{dash\_in\_host} & Creative & Presence of '-' in host; often used in brand impersonation domains. \\
\texttt{dot\_count\_host} & Creative & Number of dots in host; indicates domain depth and possible spoofing complexity. \\
\texttt{slash\_count\_path} & Creative & Count of path separators; deeper paths may correlate with deceptive navigation. \\
\texttt{percent\_encoding} & Creative & Binary indicator for URL percent encoding; sometimes used to hide suspicious strings. \\
\texttt{suspicious\_words} & Creative & Count of suspicious security/payment words (verify, update, login, secure, account). \\
\texttt{has\_https} & Creative & HTTPS usage indicator; useful but not sufficient since phishing can also use TLS. \\
\bottomrule
\end{tabular}
\end{table*}

\subsection{LaTeX Build Procedure}

The IEEE report is compiled with a standard bibliography pipeline:
\begin{lstlisting}[basicstyle=\ttfamily\footnotesize,frame=single]
cd answer/report
pdflatex -interaction=nonstopmode main.tex
bibtex main
pdflatex -interaction=nonstopmode main.tex
pdflatex -interaction=nonstopmode main.tex
pdfinfo main.pdf
\end{lstlisting}

This generates the final PDF at:
\begin{center}
\texttt{answer/report/main.pdf}
\end{center}

\subsection{Validation Checklist}

The final submission was validated against these criteria:
\begin{enumerate}
    \item Full notebook execution completes with zero cell errors.
    \item Real datasets are loaded from \texttt{datasets/} paths without fallback demo mode.
    \item All expected figure files exist and are non-empty.
    \item IEEE LaTeX report compiles successfully with bibliography.
    \item Report includes all required assignment analyses for both questions.
\end{enumerate}

\subsection{Submission-Readiness Notes}

The current artifact set is designed to be grader-friendly:
\begin{itemize}
    \item Notebook contains executable code and generated outputs.
    \item Report references only local generated figures.
    \item Paths are deterministic and independent of notebook launch location.
    \item Dependencies are explicit and minimal.
\end{itemize}

Together, these properties reduce ambiguity during evaluation and make the submission resilient to environment differences.
